% GENERATED FILE. Edit canonical lessons, then run npm run generate:books.
% canonical-source-hash: fnv1a64:d8537dcab3a05309
% canonical-lessons: AR-C549-ulliyya, AR-C549-jaras, AR-C549-miqbad, AR-C549-mizlaj, AR-C549-balua

\chapter{The attic and the door}
\label{ch:ar-the-attic-and-the-door}
\hlchaptermodality{549}

\begin{chapteropening}
\textbf{By the end of this chapter:} \emph{I can name an attic, a bell, a handle, a bolt and a drain.}

Complete the last lesson of chapter 549: I can name an attic, a bell, a handle, a bolt and a drain.
\end{chapteropening}

\section[\texorpdfstring{\ar{علية} (ʿulliyya)}{علية (ʿulliyya)}]{\ar{علية} (ʿulliyya) — an attic, a loft}

\label{lesson:AR-C549-ulliyya}



\begin{quote}
Before the new one: say the Arabic for to compete with, then the Arabic for to punch.
\end{quote}

\subsection*{You'll want to know: \ar{علية}}
\textbf{\ar{علية}} — \emph{ʿulliyya} — \textquotedblleft{}an attic, a loft\textquotedblright{}.

\begin{quote}
\ar{وضعنا} \ar{الصناديق} \ar{القديمة} \ar{في} \ar{العلية}.
\end{quote}

\begin{quote}
\emph{waḍaʿnā aṣ-ṣanādīq al-qadīma fī al-ʿulliyya.}
\end{quote}

\begin{quote}
We put the old boxes in the attic.
\end{quote}

\begin{cousinweb}
\textbf{\ar{علية}} (\emph{ʿulliyya}) is built on the native Arabic root \emph{ʿ-l-w}, \textquotedblleft{}to be high, to rise\textquotedblright{} (Lane). The ʿulliyya is the room at the top.
\end{cousinweb}

\begin{grammarlens}[title={What you've built}]
One more word for describing what you lived and felt.
\end{grammarlens}

\subsection*{Your turn}
\begin{itemize}
\raggedright
  \item \emph{Say it:} \emph{ʿulliyya}
  \item \emph{Say it:} \emph{ʿulliyya}, once more
  \item \emph{Say it:} \emph{waḍaʿnā aṣ-ṣanādīq al-qadīma fī al-ʿulliyya.} — We put the old boxes in the attic.
  \item \emph{Recall:} say the Arabic for to compete with, then the Arabic for to punch, then say \emph{ʿulliyya} again
\end{itemize}

\begin{tcolorbox}[breakable,colback=teal!4,colframe=teal!35!black,title={Before you move on}]
What does \ar{علية} mean? (\textquotedblleft{}An attic, a loft\textquotedblright{}.) Say it once more.
\end{tcolorbox}
\section[\texorpdfstring{\ar{جرس} (jaras)}{جرس (jaras)}]{\ar{جرس} (jaras) — a bell}

\label{lesson:AR-C549-jaras}



\begin{quote}
Before the new one: say the Arabic for to punch, then the Arabic for an attic.
\end{quote}

\subsection*{You'll want to know: \ar{جرس}}
\textbf{\ar{جرس}} — \emph{jaras} — \textquotedblleft{}a bell\textquotedblright{}.

\begin{quote}
\ar{رن} \ar{جرس} \ar{الباب} \ar{في} \ar{الصباح}.
\end{quote}

\begin{quote}
\emph{ranna jaras al-bāb fī aṣ-ṣabāḥ.}
\end{quote}

\begin{quote}
The doorbell rang in the morning.
\end{quote}

Plural: \ar{أجراس}, \emph{ajrās}.

\ar{جرس} \ar{الباب} is a doorbell.

\begin{cousinweb}
\textbf{\ar{جرس}} (\emph{jaras}) is built on the native Arabic root \emph{j-r-s}, \textquotedblleft{}a sound, a low voice\textquotedblright{} (Lane). Lane gives jaras as a sound, and a bell that makes it.
\end{cousinweb}

\begin{grammarlens}[title={What you've built}]
One more word for describing what you lived and felt.
\end{grammarlens}

\subsection*{Your turn}
\begin{itemize}
\raggedright
  \item \emph{Say it:} \emph{jaras}
  \item \emph{Say it:} \emph{jaras}, once more
  \item \emph{Say it:} \emph{ranna jaras al-bāb fī aṣ-ṣabāḥ.} — The doorbell rang in the morning.
  \item \emph{Recall:} say the Arabic for to punch, then the Arabic for an attic, then say \emph{jaras} again
\end{itemize}

\begin{tcolorbox}[breakable,colback=teal!4,colframe=teal!35!black,title={Before you move on}]
What does \ar{جرس} mean? (\textquotedblleft{}A bell\textquotedblright{}.) Say it once more.
\end{tcolorbox}
\section[\texorpdfstring{\ar{مقبض} (miqbaḍ)}{مقبض (miqbaḍ)}]{\ar{مقبض} (miqbaḍ) — a handle}

\label{lesson:AR-C549-miqbad}



\begin{quote}
Before the new one: say the Arabic for an attic, then the Arabic for a bell.
\end{quote}

\subsection*{You'll want to know: \ar{مقبض}}
\textbf{\ar{مقبض}} — \emph{miqbaḍ} — \textquotedblleft{}a handle\textquotedblright{}.

\begin{quote}
\ar{انكسر} \ar{مقبض} \ar{الباب}.
\end{quote}

\begin{quote}
\emph{inkasara miqbaḍ al-bāb.}
\end{quote}

\begin{quote}
The door handle broke.
\end{quote}

\begin{cousinweb}
\textbf{\ar{مقبض}} (\emph{miqbaḍ}) is built on the native Arabic root \emph{q-b-ḍ}, \textquotedblleft{}to grasp with the hand\textquotedblright{} (Lane). A miqbaḍ is the place where you grasp.
\end{cousinweb}

\begin{grammarlens}[title={What you've built}]
One more word for describing what you lived and felt.
\end{grammarlens}

\subsection*{Your turn}
\begin{itemize}
\raggedright
  \item \emph{Say it:} \emph{miqbaḍ}
  \item \emph{Say it:} \emph{miqbaḍ}, once more
  \item \emph{Say it:} \emph{inkasara miqbaḍ al-bāb.} — The door handle broke.
  \item \emph{Recall:} say the Arabic for an attic, then the Arabic for a bell, then say \emph{miqbaḍ} again
\end{itemize}

\begin{tcolorbox}[breakable,colback=teal!4,colframe=teal!35!black,title={Before you move on}]
What does \ar{مقبض} mean? (\textquotedblleft{}A handle\textquotedblright{}.) Say it once more.
\end{tcolorbox}
\section[\texorpdfstring{\ar{مزلاج} (mizlāj)}{مزلاج (mizlāj)}]{\ar{مزلاج} (mizlāj) — a bolt, a latch}

\label{lesson:AR-C549-mizlaj}



\begin{quote}
Before the new one: say the Arabic for a bell, then the Arabic for a handle.
\end{quote}

\subsection*{You'll want to know: \ar{مزلاج}}
\textbf{\ar{مزلاج}} — \emph{mizlāj} — \textquotedblleft{}a bolt, a latch\textquotedblright{}.

\begin{quote}
\ar{أغلقت} \ar{الباب} \ar{بالمزلاج}.
\end{quote}

\begin{quote}
\emph{aghlaqtu al-bāb bi-l-mizlāj.}
\end{quote}

\begin{quote}
I shut the door with the bolt.
\end{quote}

\begin{cousinweb}
\textbf{\ar{مزلاج}} (\emph{mizlāj}) is built on the native Arabic root \emph{z-l-j}, \textquotedblleft{}to slide, to slip\textquotedblright{} (Lane). Lane gives azlaja al-bāb as to make a door fast with the mizlāj. The pattern mifʿāl names a tool.
\end{cousinweb}

\begin{grammarlens}[title={What you've built}]
One more word for describing what you lived and felt.
\end{grammarlens}

\subsection*{Your turn}
\begin{itemize}
\raggedright
  \item \emph{Say it:} \emph{mizlāj}
  \item \emph{Say it:} \emph{mizlāj}, once more
  \item \emph{Say it:} \emph{aghlaqtu al-bāb bi-l-mizlāj.} — I shut the door with the bolt.
  \item \emph{Recall:} say the Arabic for a bell, then the Arabic for a handle, then say \emph{mizlāj} again
\end{itemize}

\begin{tcolorbox}[breakable,colback=teal!4,colframe=teal!35!black,title={Before you move on}]
What does \ar{مزلاج} mean? (\textquotedblleft{}A bolt, a latch\textquotedblright{}.) Say it once more.
\end{tcolorbox}
\section[\texorpdfstring{\ar{بالوعة} (bālūʿa)}{بالوعة (bālūʿa)}]{\ar{بالوعة} (bālūʿa) — a drain}

\label{lesson:AR-C549-balua}



\begin{quote}
Before the new one: say the Arabic for a handle, then the Arabic for a bolt.
\end{quote}

\subsection*{You'll want to know: \ar{بالوعة}}
\textbf{\ar{بالوعة}} — \emph{bālūʿa} — \textquotedblleft{}a drain\textquotedblright{}.

\begin{quote}
\ar{سدت} \ar{الأوراق} \ar{البالوعة} \ar{في} \ar{الشارع}.
\end{quote}

\begin{quote}
\emph{saddat al-awrāq al-bālūʿa fī ash-shāriʿ.}
\end{quote}

\begin{quote}
Leaves blocked the drain in the street.
\end{quote}

\begin{cousinweb}
\textbf{\ar{بالوعة}} (\emph{bālūʿa}) is built on the native Arabic root \emph{b-l-ʿ}, \textquotedblleft{}to swallow\textquotedblright{} (Lane). A bālūʿa swallows the water. Ibtalaʿa, to swallow, which you know, is from the same root.
\end{cousinweb}

\begin{grammarlens}[title={What you've built}]
One more word for describing what you lived and felt.
\end{grammarlens}

\subsection*{Your turn}
\begin{itemize}
\raggedright
  \item \emph{Say it:} \emph{bālūʿa}
  \item \emph{Say it:} \emph{bālūʿa}, once more
  \item \emph{Say it:} \emph{saddat al-awrāq al-bālūʿa fī ash-shāriʿ.} — Leaves blocked the drain in the street.
  \item \emph{Recall:} say the Arabic for a handle, then the Arabic for a bolt, then say \emph{bālūʿa} again
\end{itemize}

\begin{tcolorbox}[breakable,colback=teal!4,colframe=teal!35!black,title={Before you move on}]
What does \ar{بالوعة} mean? (\textquotedblleft{}A drain\textquotedblright{}.) Say it once more.
\end{tcolorbox}
